\section{Computing repeated cohort summaries}
\label{sec:performance}

The accounting definition in \eqref{eq:cohort-productivity} also determines a small computational problem. An analyst may compare many acceptance horizons on the same observed cohort while keeping the follow-up duration $\Delta$ fixed. Repeatedly scanning every record gives a simple reference implementation. Sorting the records once and storing cumulative totals replaces each subsequent scan with two binary searches. The optimisation preserves the reported quantities while changing their computational cost. Its runtime measurements do not establish an effect of AI assistance on human labour.

\subsection{A discrete input contract}

The index receives two streams with different accounting roles. Let $A=((a_i,w_i))_{i=1}^{n}$ contain one accepted-output record per accepted requirement. Each record gives an acceptance age $a_i$ measured from that requirement's enrolment and a retained weight $w_i$ assessed at the common follow-up $\Delta$. Let $B=((b_j,v_j))_{j=1}^{m}$ contain task-attributable labour contributions, including contributions from unfinished, rejected and abandoned tasks. Each $b_j$ is the elapsed age from the associated task's enrolment to its contribution's integration boundary. Shared work is attributed before these records are formed.

Ages use a declared common integer time unit. All ages and weights are nonnegative \texttt{u64} integers, with output weights in millionths of a reference unit and labour weights in millionths of a person-hour. The two weight types are summed separately.

For nonnegative \texttt{u64} horizons $\tau$ and $\Delta$ in that time unit, the computational targets are
\begin{equation}
\label{eq:discrete-cohort-totals}
 Q_{\mathrm{d}}(\tau)=\sum_{i=1}^{n} w_i\ind\{a_i\leq\tau\},
 \qquad
 H_{\mathrm{d}}(\tau,\Delta)=\sum_{j=1}^{m}v_j\ind\{b_j\leq\tau+\Delta\}.
\end{equation}
The result is the exact integer pair $(Q_{\mathrm{d}},H_{\mathrm{d}})$. Their common scale cancels in the productivity ratio, which is defined only when $H_{\mathrm{d}}>0$. A zero denominator remains an explicit zero in the pair. The Rust convenience method \texttt{ratio()} returns \texttt{None} in that case. Its floating-point result when the denominator is positive is a display approximation, not the canonical stored value.

To relate these records to the continuous integral in \eqref{eq:labour}, the observation and integration rule must supply the correct contributions at each queried boundary. Cutoffs aligned with already integrated bins are one example. Assigning a whole contribution to the end of a sampled interval or time bin instead approximates an interior cutoff.

The measurement protocol must also resolve missing labour, unobserved retention, duplicated tasks and scope changes, which the index cannot diagnose. Retained weights apply to the particular $\Delta$ at which they were assessed. A change to $\Delta$, the cohort or its assessed weights requires corresponding input data and a rebuilt index. The next result states what exact bin masses can establish at an interior cutoff.

\subsection{Bounding labour at a bin boundary}
\label{sec:binning-bounds}

Exact bin masses can bound the continuous labour target even when the distribution of work within each bin is unknown. Suppose a fixed cohort's labour is partitioned into complete, nonduplicated contributions on task-relative age intervals $[\ell_j,u_j)$, with $0\leq\ell_j<u_j$. Attribute shared work once. Each contribution arises from a nonnegative allocation integral with absolutely continuous cumulative labour and known total mass $v_j$ in person-hours. Coverage and attribution are premises of the bound, not facts established by the records alone. Ages and cutoffs use a common declared time unit, and masses must remain exact after conversion. Rounded weights require separate error bounds.

\begin{proposition}[A continuous labour envelope]
\label{prop:binning-bounds}
Under these premises, set $c=\tau+\Delta$ and $H(c)=\sum_i H_i(c)$. Then
\begin{equation}
\label{eq:binning-envelope}
 H_-(c):=\sum_jv_j\ind\{u_j\leq c\}
 \;\leq H(c)\leq\;
 H_+(c):=\sum_jv_j\ind\{\ell_j<c\}.
\end{equation}
The gap is $H_+(c)-H_-(c)=\sum_{j:\,\ell_j<c<u_j}v_j$. If $H_-(c)>0$ and retained output $Q=\sum_iQ_i(\tau,\Delta)$ is known, then
\begin{equation}
\label{eq:binning-ratio}
 \frac{Q}{H_+(c)}\leq P_H(\tau,\Delta)\leq\frac{Q}{H_-(c)}.
\end{equation}
\end{proposition}
\begin{proof}
A bin ending by $c$ contributes its full mass. A bin starting at or after $c$ contributes zero; absolute continuity excludes endpoint atoms. Every crossing bin contributes between zero and its mass. Summing gives \eqref{eq:binning-envelope} and the gap. Nonnegative $Q$ and positive bounding denominators give \eqref{eq:binning-ratio} by division.
\end{proof}

When $H_-=0$, the bounds need not establish a positive denominator or a finite upper productivity bound. If output instead has justified bounds $Q_-\leq Q\leq Q_+$, positive $H_-$ gives $Q_-/H_+\leq P_H\leq Q_+/H_-$. These are deterministic binning bounds, not causal confidence intervals. They do not guarantee sampling accuracy, remove nonresponse or establish that Horvitz--Thompson estimates equal complete-population bin masses.

\begin{example}[An interior cutoff]
\label{ex:binning-cutoff}
One person's bin $[0,2)$ contains one person-hour. At $c=1$, work entirely in the first hour gives labour 1, whereas work entirely in the second gives 0. Both respect the allocation constraint. Uniform interpolation gives $1/2$, but the known bin mass alone supports the interval $[0,1]$.
\end{example}


\subsection{Construction, queries and correctness}

The sorted prefix-index constructor, \autoref{alg:build-horizon-index}, is a deterministic preprocessing algorithm for the two fixed record streams. The order of records with equal ages is immaterial: every such record is included or excluded together. Arrays use one-based record positions and a zero prefix. \Call{CheckedAdd128}{} widens the weight and either returns the exact sum or reports \texttt{TotalOverflow}. No wrapped or partial index is returned. The Rust implementation copies the input streams before sorting, so construction preserves the caller's records.

\begingroup
% A hanging contract aligns continuation lines after its semantic label.
\newlength{\horizoncontractwidth}
\newcommand{\HorizonContract}[2]{%
  \Statex\parbox[t]{\linewidth}{%
    \settowidth{\horizoncontractwidth}{\textbf{#1:}\ }%
    \hangindent=\horizoncontractwidth\hangafter=1
    \textbf{#1:}\ #2}}
\algrenewcommand{\Require}[1]{\HorizonContract{Require}{#1}}
\algrenewcommand{\Ensure}[1]{\HorizonContract{Ensure}{#1}}

\begin{algorithm}[tb]
\caption[Build a cohort horizon index]{Build exact cumulative totals for a fixed cohort and follow-up}
\label{alg:build-horizon-index}
\begin{algorithmic}[1]
\Require{Accepted-output records $A$ and labour records $B$ under the stated contract.}
\Ensure{Sorted keys and exact prefix totals, or \texttt{TotalOverflow}.}
\ForAll{stream names $S$ in the ordered list $(A,B)$}
  \State $X_S\gets\Call{SortCopyByAge}{S}$
  \State $K_S\gets$ age array of $X_S$; $P_S[0]\gets0$
  \For{$j\gets1$ \textbf{to} $|X_S|$}
    \State $z\gets\Call{CheckedAdd128}{P_S[j-1],X_S[j].\mathrm{weight}}$
    \If{$z=\mathrm{overflow}$}
      \State \Return \texttt{TotalOverflow}
    \EndIf
    \State $P_S[j]\gets z$
  \EndFor
\EndFor
\State \Return $I=(K_A,P_A,K_B,P_B)$
\end{algorithmic}
\end{algorithm}

The batch horizon evaluator, \autoref{alg:query-horizon-index}, is a read-only query algorithm on this index. \Call{UpperBound}{$K,x$} returns the number of sorted keys at most $x$, including every duplicate at the boundary. \Call{CheckedAdd64}{} reports \texttt{HorizonOverflow} if $\tau+\Delta$ cannot be represented. The entire batch is checked before any result is returned. Query order is preserved, so horizons need not arrive sorted.

\begin{algorithm}[tb]
\caption[Query a cohort horizon index]{Return exact output and labour totals for a batch of horizons}
\label{alg:query-horizon-index}
\begin{algorithmic}[1]
\Require{Index $I$; horizons $\tau_1,\ldots,\tau_q$; the fixed follow-up $\Delta$.}
\Ensure{Ordered exact total pairs, or \texttt{HorizonOverflow}; no division by zero.}
\For{$r\gets1$ \textbf{to} $q$}
  \State $c_r\gets\Call{CheckedAdd64}{\tau_r,\Delta}$
  \If{$c_r=\mathrm{overflow}$}
    \State \Return \texttt{HorizonOverflow}
  \EndIf
\EndFor
\State $O\gets$ empty result array
\For{$r\gets1$ \textbf{to} $q$}
  \State $k_A\gets\Call{UpperBound}{K_A,\tau_r}$
  \State $k_B\gets\Call{UpperBound}{K_B,c_r}$
  \State Append $(P_A[k_A],P_B[k_B])$ to $O$
\EndFor
\State \Return $O$
\end{algorithmic}
\end{algorithm}
\endgroup

\begin{proposition}[Agreement with a full scan]
\label{prop:horizon-index}
If index construction and horizon addition succeed, each returned pair in \autoref{alg:query-horizon-index} equals \eqref{eq:discrete-cohort-totals}. This includes empty streams, repeated ages, zero weights and zero labour totals.
\end{proposition}
\begin{proof}
After construction has processed $j$ records of a sorted stream, its loop invariant is $P_S[j]=\sum_{\ell=1}^{j}X_S[\ell].\mathrm{weight}$. It holds at $j=0$ and checked addition preserves it. Sorting preserves the multiset of records. For a cutoff $x$, upper-bound binary search maintains a lower prefix whose keys are at most $x$ and an upper suffix whose keys exceed $x$. Its remaining interval strictly decreases until its bounds coincide at $k$. Thus exactly positions $1,\ldots,k$ meet the inclusive cutoff. The corresponding prefix total is the same sum a full scan selects. Apply this argument at $x=\tau$ for output and $x=\tau+\Delta$ for labour. The empty prefix has total zero, so the boundary cases follow without division.
\end{proof}

The proof concerns these discrete sums and the published operations. It is an analytical proof, not an extension of the Lean proof inventory in Appendix~\ref{app:coverage}. The Rust tests check implementation agreement with a separate scanning path. Successful tests do not establish completeness or reliability of the underlying observations.

With $n$ acceptance records, $m$ labour records and $q$ queries, repeated scanning takes $O(q(1+n+m))$ time. Construction takes $O(n\log(n+1)+m\log(m+1))$ time and $O(n+m)$ additional storage. Indexed queries take $O(q(1+\log(n+1)+\log(m+1)))$ time and $O(q)$ result storage. These bounds count fixed-width operations and assume available memory. Preprocessing can outweigh the saving for a small batch, and an index cannot be reused after its input changes without maintaining or rebuilding its prefix totals.

\subsection{A reproducible performance comparison}

The supplement's dependency-free Rust program compares the scanning and indexed implementations on the same generated records and horizon batches. It retains the exact totals for equality checks and times computation separately from data generation. The benchmark reports index construction as well as query time: a fast query alone does not establish that building the index was worthwhile. The fixed follow-up and synthetic inputs isolate the scan-versus-index choice within the Rust program. They do not represent a company workload. For further reading on data science with Rust, see Van Der Post~\cite{vanderpost2024datascience}. The runtime claims here rely on the accompanying measurements.

The workload contains 100,000 generated accepted-output records and 400,000 generated labour increments. It has no task roster or task-to-increment associations. Ages and horizons use hours, with $\Delta=720$ and queried horizons from 0 to 1,440. Three fixed seeds (17, 41 and 93) define the datasets and horizon sequences. For each seed and each $q\in\{1,10,100,1000\}$, five fresh processes each run one paired comparison. Each process first performs an unmeasured full scan followed by an unmeasured indexed run. Measured method order then alternates across observations. These are timings on warmed data and code, not cold-start measurements. Both implementations return identical integer pairs in every warmup and measured comparison. All 60 measured pairs are retained.

The build uses Rust 1.85.1, optimisation level 3, one code-generation unit, no link-time optimisation and enabled overflow checks, on a shared Linux host reporting an Intel Xeon Platinum 8573C processor. Dataset generation, startup, warmup, output verification and result I/O are outside both timed regions. Indexed total time includes input copies, sorting, prefix construction, queries and disposal of the index. The environment record retains the compiler details, source hashes and timing boundaries. Host load is uncontrolled. The 15 pairs per query count are repeated measurements of three synthetic datasets, not independent samples of production work.

\begin{table}[tb]
\centering
\caption{Measured cohort-query times in milliseconds: medians of 15 paired observations at each query count. Indexed total is timed separately and includes index disposal. Component medians need not sum to it. Reduction is $100(1-\text{indexed median}/\text{scan median})$, so a negative value denotes a regression.}
\label{tab:horizon-performance}
\small
% Generated by performance/summarise.py; all times in milliseconds.
\begin{tabular}{rrrrrr}
\toprule
$q$ & Scan & Build & Query & Indexed total & Reduction \\
\midrule
1 & 2.889 & 20.348 & 0.002 & 21.016 & -627.4\% \\
10 & 28.003 & 20.206 & 0.008 & 20.709 & 26.0\% \\
100 & 263.415 & 22.348 & 0.082 & 23.046 & 91.3\% \\
1000 & 2787.159 & 21.803 & 0.508 & 23.051 & 99.2\% \\
\bottomrule
\end{tabular}

\end{table}

Table~\ref{tab:horizon-performance} retains the adverse small-batch result. For one query, indexing increases median elapsed time from 2.889 to 21.016~ms. At ten queries, the median decreases by 26.0\%, but individual paired reductions range from $-16.4\%$ to $55.1\%$: indexing does not win every observed comparison. At 100 and 1,000 queries, every measured pair favours the index. For 1,000 queries, the median falls from 2,787.159 to 23.051~ms, a 99.2\% reduction including construction and disposal. Figure~\ref{fig:horizon-performance} shows the observed ranges as well as the medians. The ranges describe variation across seeds and repeated host runs. They are not confidence intervals. The tested query counts do not identify an exact break-even point.

\begin{figure}[tb]
\centering
% Include-ready figure: one picture, no float/caption or document preamble.
% Requires pgfplots (compat=1.18); compile from paper/article, as main.tex does.
% Target: exactly 152.4 mm (6 in) wide; text >=8 pt, strokes >=0.4 pt.
% Colour and greyscale: dashed squares = repeated scans; solid diamonds = index.
% Data and finding title are generated from measured timings by prepare.py.
\begin{tikzpicture}
  \definecolor{horizonInk}{RGB}{32,48,56}
  \definecolor{horizonTeal}{RGB}{0,124,131}
  \definecolor{horizonGrey}{RGB}{110,117,120}
  % Generated by performance/figure/prepare.py from retained raw timings.
\def\HorizonFirstWin{100}
\def\HorizonFinding{Indexing repays its build cost by 100 queries}
\def\HorizonOneRatio{0.12}
\def\HorizonManyRatio{134.5}
\def\HorizonYmin{0.705847}
\def\HorizonYmax{5185.56}
\def\HorizonYticks{1,10,100,1000}
\def\HorizonRmin{0.01}
\def\HorizonRmax{1000}
\def\HorizonRticks{0.01,0.1,1,10,100,1000}
\def\HorizonScanLast{2787.15866}
\def\HorizonIndexedLast{23.050666}

  \path[use as bounding box] (0,0) rectangle (152.4mm,85mm);
  \node[anchor=base west,inner sep=0pt,text=horizonInk,
        font=\sffamily\bfseries\fontsize{10}{12}\selectfont]
    at (17mm,80mm) {\HorizonFinding};
  \node[anchor=base west,inner sep=0pt,text=horizonGrey,
        font=\sffamily\fontsize{8}{10}\selectfont]
    at (17mm,74.5mm) {Observed machine timings on generated records};
  \begin{axis}[
      at={(17mm,16mm)},anchor=south west,
      width=100mm,height=51mm,scale only axis,
      xmode=log,log basis x=10,ymode=log,log basis y=10,
      xmin=1,xmax=1000,ymin=\HorizonYmin,ymax=\HorizonYmax,
      xtick={1,10,100,1000},xticklabels={1,10,100,{1,000}},
      ytick/.expanded=\HorizonYticks,
      yticklabels/.expanded=\HorizonYticks,
      xlabel={Queries over the same cohort (log scale)},
      ylabel={Total elapsed time (ms, log scale)},
      tick label style={font=\sffamily\fontsize{8}{10}\selectfont,text=horizonInk},
      label style={font=\sffamily\fontsize{8}{10}\selectfont,text=horizonInk},
      xlabel style={yshift=1mm},ylabel style={yshift=-1mm},
      axis x line*=bottom,axis y line*=left,axis line style={horizonGrey,line width=.4pt},
      tick style={horizonGrey,line width=.4pt},tick align=outside,
      minor tick num=0,minor x tick num=0,minor y tick num=0,
      ymajorgrids=true,grid style={black!16,line width=.4pt},
      clip=false,
    ]
    \addplot+[horizonInk,dashed,line width=.85pt,mark=square*,mark size=2.2pt,
      mark options={solid,fill=horizonInk,draw=horizonInk},
      error bars/.cd,y dir=both,y explicit,
      error bar style={line width=.5pt},error mark options={line width=.5pt,mark size=2pt}]
      table[x=q,y=scan,y error minus=scan_minus,y error plus=scan_plus,col sep=comma]
      {../performance/figure/horizon-data.csv};
    \addplot+[horizonTeal,solid,line width=1pt,mark=diamond*,mark size=2.8pt,
      mark options={solid,fill=horizonTeal,draw=horizonTeal},
      error bars/.cd,y dir=both,y explicit,
      error bar style={line width=.5pt},error mark options={line width=.5pt,mark size=2pt}]
      table[x=q,y=indexed,y error minus=indexed_minus,y error plus=indexed_plus,col sep=comma]
      {../performance/figure/horizon-data.csv};
    \node[anchor=west,inner sep=0pt,text=horizonInk,
      font=\sffamily\fontsize{8}{10}\selectfont,xshift=3mm]
      at (axis cs:1000,\HorizonScanLast) {Scan each query};
    \node[anchor=west,inner sep=0pt,text=horizonTeal,
      font=\sffamily\fontsize{8}{10}\selectfont,xshift=3mm,align=left]
      at (axis cs:1000,\HorizonIndexedLast) {Build index\,+\,queries};
  \end{axis}
  \node[anchor=base west,inner sep=0pt,text=horizonGrey,
        font=\sffamily\fontsize{8}{10}\selectfont]
    at (17mm,1.2mm) {Medians and observed ranges; build and disposal included.};
\end{tikzpicture}%

\caption[Measured cohort-query times]{Observed elapsed times for repeated cohort-horizon queries on generated arrays of 100,000 accepted-output records and 400,000 labour increments. Markers are medians. Whiskers show the minimum and maximum across 15 paired runs (three seeds, five fresh-process repetitions per seed). They are not confidence intervals. Indexed totals include construction, queries and disposal. All one-query pairs favoured scanning and all 100-query pairs favoured indexing. Ten-query pairs varied in direction. Only four query counts were timed, so the connecting lines do not locate break-even. This result is confined to the generated workload and execution environment and does not estimate human labour or an effect of an AI intervention.
}
\label{fig:horizon-performance}
\end{figure}

The observed saving is consistent with amortising preprocessing over repeated scans. It supports the index for this fixed workload and query pattern, subject to its storage and update costs. It does not establish performance on larger-than-memory datasets, continuously changing ledgers, different hardware or different query distributions. Raw paired timings, the generator, exact-output checks and the summary script are retained so that another workload can test those boundaries.
